\subsection{\texttt{TRSQRT}}
\noindent\textbf{Bundle encoding.} UOP entry 56, opcode \texttt{0x1055}. Bundle: \texttt{B.OP + B.ARG + B.IOTI}. \texttt{B.OP.Opcode}=\texttt{0x1055}. \texttt{B.IOTI}: selector=\texttt{111 last}, \texttt{SrcTile0}=\texttt{src}, \texttt{SrcTile1}=\texttt{absent}, \texttt{DstTile}=\texttt{dst}, \texttt{SizeCode}=profile tile size.\par\smallskip

\textit{Semantics.} Computes the element-wise reciprocal square root of \texttt{src} as \texttt{1.0 / sqrt(src[r, c])} and writes the result to \texttt{dst}. The microcode body achieves this by computing \texttt{sqrt(src)} then dividing 1.0 by the result. Note: for f16, the microcode body uses f16 for the broadcast 1.0; for f32, the microcode body uses f32. Mixing element formats across operands MUST be performed by an explicit \texttt{TCVT} step.

\textit{Domain errors.} When \texttt{src[r, c] < 0}, the result is NaN. When \texttt{src[r, c]} is +Inf, the result is zero. When \texttt{src[r, c]} is +0, the result is +Inf. When \texttt{src[r, c]} is NaN, the result is NaN. When \texttt{src[r, c]} is negative, the result is NaN.\par\smallskip
{\small
\paragraph{PTO ISA description.}
Elementwise reciprocal square root.
\paragraph{Example.}
Tile elements [1, 4, 9] give [1.0, 0.5, 0.333...].
\par\smallskip
\paragraph{PTO ISA operation.}
For each element \texttt{(i, j)} in the valid region:

\[
\mathrm{dst}_{i,j} = \frac{1}{\sqrt{\mathrm{src}_{i,j}}}
\]
}\par\smallskip
\paragraph{Notes.}
1/0 = Inf. 1/sqrt(0) = Inf. 1/sqrt(Inf) = 0. Negative input produces NaN. Accuracy is implementation-defined; hardware implementations typically use a Newton-Raphson refinement.
\paragraph{Microcode site table.}
{\scriptsize
\begin{longtable}{@{}R{0.055\linewidth}L{0.23\linewidth}L{0.31\linewidth}L{0.22\linewidth}@{}}
\toprule
Seq & Micro instruction & WO[63:32] fields & WM[31:0] fields \\
\midrule
0 & \texttt{u\_vec.vlds} & \texttt{opc=0x17, x=0, fl=mem, seq=0, kind=b} & \texttt{entry=56, cnt=2, res=vec, var=0} \\
1 & \texttt{u\_vec.vbr} & \texttt{opc=0x0B, x=0, fl=alu, seq=1, kind=c} & \texttt{entry=56, cnt=0, res=vec, var=0} \\
2 & \texttt{u\_vec.vsqrt} & \texttt{opc=0x29, x=0, fl=alu, seq=2, kind=u} & \texttt{entry=56, cnt=1, res=vec, var=0} \\
3 & \texttt{u\_vec.vdiv} & \texttt{opc=0x12, x=0, fl=alu, seq=3, kind=b} & \texttt{entry=56, cnt=2, res=vec, var=0} \\
4 & \texttt{u\_vec.vsts} & \texttt{opc=0x2A, x=0, fl=mem, seq=4, kind=b} & \texttt{entry=56, cnt=2, res=vec, var=0} \\
\bottomrule
\end{longtable}
}
\paragraph{Microcode body DSL.}
\begin{lstlisting}[style=ptocode]
def uop_trsqrt(src, dst):
    dtype = dst.element_type
    valid_rows, valid_cols = dst.valid_shape

    for row in range(0, valid_rows, 1):
        remained = valid_cols
        for col in range(0, valid_cols, pto.get_lanes(dtype)):
            mask, remained = pto.make_mask(dtype, remained)
            vinput = pto.vlds(src[row, col:])
            if pto.constexpr(dtype == pto.f16):
                one_scalar = pto.f16(1.0)
            else:
                one_scalar = pto.f32(1.0)
            one = pto.vbr(one_scalar)
            sqrt_result = pto.vsqrt(vinput, mask)
            result = pto.vdiv(one, sqrt_result, mask)
            pto.vsts(result, dst[row, col:], mask)
    return
\end{lstlisting}
\Needspace{0.34\textheight}
